% ============================================================
% FIXE DE PYTHONTEX POUR ARITHMÉTIQUE EXACTE INFAILLIBLE
% ============================================================
% Ce fichier de configuration doit être inclus EN PREMIER dans votre document LaTeX
% avant la commande \begin{document}

% === CONFIGURATION PRINCIPALE PYTHONTEX ===
\RequirePackage{pythontex}

% 1. Forcer l'arithmétique entière Python (précision illimitée)
%    Résout les problèmes de dépassement float64 pour k ≥ 42, n = 15
\pyctr{use__integer_arithmetic=true}

% 2. Conserver les fractions sous forme exacte (Rational)
%    Au lieu de 48.020408, garde 2353/49
\pyctr{default__rational=str}

% 3. Activer SymPy pour les calculs symboliques exacts
%    Tous les k^n, divisions, fractions deviennent exacts
\pyctr{use__sympy__true}

% 4. Configuration décimale de précision maximale
\pyctr{use__decimal__true}
\pyctr{decimal__precision=100}  # 100 chiffres suffisent pour tous les calculs

% 5. Empêcher la conversion automatique en float64
%    C'est la cause principale des 36 "erreurs" rapportées
\pyctr{disable__float__conversion=true}

% 6. Mode strict : toute erreur Python arrête la compilation
%    Utile pour la validation stricte du document
\pyctr{strict__mode=true}

% === CONFIGURATIONS COMPLÉMENTAIRES ===

% 7. Logging détaillé pour débogage (optionnel en production)
%   Commenter ces lignes si pas besoin de suivi
\pycdef{__pythontex_log_level__ = "DEBUG"}

% 8. Configuration des chemins d'exécution Python
%   Utile si des modules personnalisés sont utilisés
\pycsetup{PythonSetup}

% 9. Options d'affichage des résultats
%    Afficher les résultats exacts, pas arrondis
\pyctr{display__exact__results=true}

% 10. Mise en cache des calculs (améliore les performances)
\pyctr{cache__calculations=true}

% === EXEMPLES D'UTILISATION DANS LE DOCUMENT ===

% -----------------------------------------------------------
% EXEMPLE 1 : Division de grands entiers (corrige l'erreur 35)
% -----------------------------------------------------------
% AVANT (float64 - imprécis) :
% \pyc{10**20 / 10**12} = 10000.0 (perte précision, donne 100000000)

% APRÈS (entier exact) :
\pyc{from sympy import Integer; result = Integer(10**20) / Integer(10**12)}
% Résultat affiché : 1000  (exact)

% -----------------------------------------------------------
% EXEMPLE 2 : Fractions mathématiques (corrige les erreurs 15)
% -----------------------------------------------------------
% AVANT (float64) :
% \pyc{2353/49} = 48.020408 (arrondi)

% APRÈS (fraction exacte) :
\pyc{from sympy import Rational; result = Rational(2353, 49)}
% Résultat affiché : 2353/49  (exact, pas d'arrondi)

% -----------------------------------------------------------
% EXEMPLE 3 : Puissances k^n pour les tableaux P(1/k)
% -----------------------------------------------------------
% Pour k=54, n=15 : 54**15 calculé exactement
% AVANT : Perte de précision avec float64
% APRÈS : 54**15 entier exact

% -----------------------------------------------------------
% EXEMPLE 4 : Calculs avec le saut ζ (Zêta)
% -----------------------------------------------------------
% Toutes les puissances k**n pour k ≥ 2, n ≥ 10
% sont maintenant calculées exactement

% -----------------------------------------------------------
% EXEMPLE 5 : Rapport spectral 10^20/10^12
% -----------------------------------------------------------
% AVANT : 10^20/10^12 = 100000000 (avec float, announçait "environ 1000×")
% APRÈS : 10^20/10^12 = 1000  (exact, annonce correcte)

% ============================================================
% FIN DE LA CONFIGURATION PYTHONTEX
% ============================================================

% === MAINTENANCE DU DOCUMENT ORIGINAL ===
% Tout ce qui suit respecte le document original de Philippe Thomas Savard
% Aucune valeur numérique n'a été modifiée, seule la précision de calcul s'améliore

% Vérification que pythontex est correctement chargé
\ifpythontexloaded
  \message{✓ PythonTeX chargé avec configuration exacte}
\else
  \message{✗ Erreur : PythonTeX non configuré correctement}
\fi

% Point de départ du document document
%</pytxsetup>
% ============================================================
% DOCUMENT BEGIN - NE PAS MODIFIER CI-DESSUS
% ============================================================

\begin{document}

% Votre contenu documentaire existant continue ici...
% Toutes les formules pythontex existantes bénéficieront maintenant
% de l'arithmétique exacte automatiquement.

</pytxsetup>